\setcounter{chapter}{8}
\section*{Capitulo 9. Derivadas y aplicaciones}
\addcontentsline{toc}{section}{Capitulo 9. Derivadas y aplicaciones}

\begin{shortanswer}{[GX-C09.S03-01]}
Tangente: \(y=-2x\). Normal: \(y=\dfrac{x}{2}+\dfrac{5}{2}\).
\end{shortanswer}

\begin{shortanswer}{[GX-C09.S03-02]}
Tangente: \(y=-2\). Normal: \(x=1\).
\end{shortanswer}

\begin{shortanswer}{C09.S03 practica autonoma}
\begin{enumerate}[label=\textbf{PX-C09.S03-\arabic*:}, leftmargin=*]
  \item Tangente: \(y=6x-4\). Normal: \(y=-\dfrac{x}{6}+\dfrac{25}{3}\).
  \item Tangente: \(y=-\dfrac{x}{2}+\dfrac{7}{2}\). Normal: \(y=2x-4\).
  \item Tangente: \(y=\dfrac{x}{4}+\dfrac{5}{4}\). Normal: \(y=-4x+14\).
  \item Tangente: \(y=x+1\). Normal: \(y=-x+1\).
  \item Tangente: \(y=x-1\). Normal: \(y=-x+1\).
  \item Tangente: \(y=0\). Normal: \(x=0\).
\end{enumerate}
\end{shortanswer}

\begin{shortanswer}{[GX-C09.S04-01]}
\(y=2x-9\).
\end{shortanswer}

\begin{shortanswer}{[GX-C09.S04-02]}
\(y=\dfrac{x}{4}+2\).
\end{shortanswer}

\begin{shortanswer}{C09.S04 practica autonoma}
\begin{enumerate}[label=\textbf{PX-C09.S04-\arabic*:}, leftmargin=*]
  \item \(y=4x-3\).
  \item \(y=12x+16\) y \(y=12x-16\).
  \item \(y=x-1\).
  \item \(y=e^2x-e^2\).
  \item \(y=2\) y \(y=-2\).
  \item \(a=3\) y la tangente es \(y=5x-1\).
\end{enumerate}
\end{shortanswer}

\begin{shortanswer}{[CH-C09-01]}
La tangente es horizontal en \(t=3\), su ecuacion es \(y=14\) y representa la altura maxima
instantanea de la trayectoria.
\end{shortanswer}
